\subsection{Homogeneous slopes}

Table~\ref{tab:mc1} compares the 2SIV estimator with the first-stage IV
estimator and with IV estimation that ignores the factors
(\code{factmax=0}), for $(N,T)\in\{(50,50),(100,100),(200,200)\}$. Size
refers to two-sided $t$ tests at the 5\% level, and the last column reports
the rejection frequency of the $J$ test at the 5\% level.

\begin{table}[t!]
\centering\scriptsize
\begin{tabular}{llrrrrrrrrrr}
\toprule
 & & \multicolumn{3}{c}{$\psi_0$} & \multicolumn{3}{c}{$\rho_1$} & \multicolumn{3}{c}{$\beta_1$} & \\
\cmidrule(lr){3-5}\cmidrule(lr){6-8}\cmidrule(lr){9-11}
$(N,T)$ & Estimator & Bias & RMSE & Size & Bias & RMSE & Size & Bias & RMSE & Size & $J$ \\
\midrule
$(50,50)$ & 2SIV & 0.005 & 0.034 & 9.1 & $-$0.003 & 0.020 & 11.1 & 0.002 & 0.026 & 10.4 & 4.9 \\
 & First-stage IV & 0.021 & 0.050 & 12.3 & $-$0.018 & 0.030 & 8.1 & 0.012 & 0.029 & 7.5 & 18.0 \\
 & IV, no factors & 0.028 & 0.036 & 41.6 & $-$0.010 & 0.016 & 30.4 & 0.354 & 0.357 & 100.0 & 21.5 \\
\addlinespace
$(100,100)$ & 2SIV & 0.001 & 0.015 & 6.5 & $-$0.001 & 0.009 & 9.6 & 0.000 & 0.012 & 7.6 & 3.5 \\
 & First-stage IV & 0.003 & 0.018 & 6.6 & $-$0.009 & 0.016 & 7.5 & 0.005 & 0.014 & 6.5 & 14.8 \\
 & IV, no factors & 0.019 & 0.023 & 45.6 & $-$0.007 & 0.010 & 40.6 & 0.360 & 0.361 & 100.0 & 54.0 \\
\addlinespace
$(200,200)$ & 2SIV & 0.000 & 0.007 & 5.8 & 0.000 & 0.004 & 5.3 & 0.000 & 0.006 & 4.5 & 3.8 \\
 & First-stage IV & 0.001 & 0.008 & 6.6 & $-$0.004 & 0.007 & 5.5 & 0.003 & 0.007 & 6.2 & 17.4 \\
 & IV, no factors & 0.012 & 0.015 & 49.2 & $-$0.004 & 0.007 & 46.1 & 0.363 & 0.363 & 100.0 & 91.1 \\
\bottomrule
\end{tabular}

\caption{Homogeneous slopes: bias, root mean squared error, and size (\%) of
5\% $t$ tests for $\psi_0=0.3$, $\rho_1=0.4$, $\beta_1=1$, and rejection
frequency (\%) of the $J$ test. 1{,}000 replications.}
\label{tab:mc1}
\end{table}

The 2SIV estimator is practically unbiased in all designs, and its RMSE
halves as $N$ and $T$ double. Its $t$ tests over-reject somewhat in the
smallest design (sizes between 9.1\% and 11.1\%), but their size is between
4.5\% and 5.8\% when $N=T=200$, and the $J$ test has size between 3.5\% and
4.9\%. The eigenvalue-ratio criterion selects the true number of factors in
the instruments in at least 99.8\% of the replications and in the residuals
in at least 97.2\%. The first-stage estimator is consistent but less
precise; its $J$ test rejects in 15\%--18\% of the replications because it
ignores the factors in the error, which illustrates why the second stage is
the default. Ignoring the factors produces a bias of about 0.36 in
$\beta_1$ that does not vanish as the sample grows, so that $t$ tests on
$\beta_1$ always reject. Reassuringly, the $J$ test detects this
misspecification with a frequency that rises from 21.5\% to 91.1\%.

\subsection{Inference on long-run impacts}

To make the long-run multipliers strongly non-linear, Table~\ref{tab:mc2}
uses $\psi_0=0.35$ and $\rho_1=0.5$, so that $1-\rho_1-\psi_0=0.15$ and the
total long-run impact of $x_1$ is $6.67$. It reports the coverage of
nominal 95\% intervals, the frequencies with which the true value lies below
or above the interval, and the median interval width. In the smallest
design, the point estimates violated the stability condition in 22 of the
1{,}000 replications, for which impacts are undefined; these replications are
excluded.

\begin{table}[t!]
\centering\scriptsize
\begin{tabular}{llrrrrrrrr}
\toprule
 & & \multicolumn{4}{c}{Delta method} & \multicolumn{4}{c}{Simulation (Krinsky--Robb)} \\
\cmidrule(lr){3-6}\cmidrule(lr){7-10}
$(N,T)$ & Impact & Cover & Below & Above & Width & Cover & Below & Above & Width \\
\midrule
$(30,30)$ & direct & 88.7 & 1.7 & 9.6 & 0.63 & 87.2 & 7.1 & 5.7 & 0.71 \\
 & indirect & 89.1 & 0.2 & 10.7 & 5.60 & 86.9 & 7.6 & 5.5 & 7.15 \\
 & total & 88.7 & 0.2 & 11.1 & 6.05 & 87.1 & 7.5 & 5.4 & 7.60 \\
\addlinespace
$(50,50)$ & direct & 92.7 & 1.6 & 5.7 & 0.40 & 90.4 & 5.6 & 4.0 & 0.41 \\
 & indirect & 90.8 & 0.5 & 8.7 & 3.76 & 88.6 & 6.4 & 5.0 & 4.15 \\
 & total & 90.8 & 0.5 & 8.7 & 4.08 & 88.3 & 6.4 & 5.3 & 4.48 \\
\addlinespace
$(100,100)$ & direct & 92.9 & 2.2 & 4.9 & 0.19 & 93.7 & 3.2 & 3.1 & 0.19 \\
 & indirect & 93.3 & 1.4 & 5.3 & 1.85 & 94.5 & 2.7 & 2.8 & 1.89 \\
 & total & 93.5 & 1.4 & 5.1 & 2.00 & 94.2 & 2.6 & 3.2 & 2.04 \\
\bottomrule
\end{tabular}

\caption{Long-run impacts of $x_1$ with $\psi_0=0.35$, $\rho_1=0.5$:
coverage (\%) of nominal 95\% intervals, frequency (\%) with which the true
value lies below the lower or above the upper limit, and median width.}
\label{tab:mc2}
\end{table}

The results give a nuanced picture that users should know about. The Delta
intervals are symmetric while the sampling distribution of the impacts is
skewed to the right; as a result almost all of their non-coverage is on one
side: the true value lies above the upper limit in 9\%--11\% of the
replications for indirect and total impacts when $N,T\le 50$, and below the
lower limit in less than 1\%. Simulation intervals spread the errors almost
evenly between the two tails. Their overall coverage, however, is not
better in small samples (86.9\%--88.6\% compared with 88.7\%--90.8\% for the Delta
method), because both methods inherit the same first-order approximation of
the variance. In the largest design both methods are close to nominal
coverage, simulation slightly more so (94.2\% compared with 93.5\% for total
impacts). The practical recommendation is therefore to report simulation
intervals when the process is persistent, because they convey the
asymmetry of the uncertainty, while recognizing that neither method is
exact in small samples.

\subsection{Mean-group estimation with aberrant units}

The data for Table~\ref{tab:mc3} have heterogeneous slopes, with
$\psi_{0i}$ and $\rho_{1i}$ uniform on $\pm0.1$ around their means and
$\beta_{\ell i}$ uniform on $\pm0.4$ around $\beta_\ell$. A share of 0\%,
5\%, or 10\% of the units is \emph{aberrant}: both of their covariate slopes
are shifted by $+3$. The target is the mean of the regular units,
$\beta_1=1$. The table reports the MG estimator, the coefficient-wise 10\%
trimmed mean and the median of the unit-specific estimates (with 499
bootstrap replications over units), and the MG estimator after removing the
units flagged on $\beta_1$ or $\beta_2$, together with the true and false
positive rates of the flags.

\begin{table}[t!]
\centering\scriptsize
{\setlength{\tabcolsep}{4pt}\begin{tabular}{llrrrrrrrrr}
\toprule
 & & \multicolumn{3}{c}{No aberrant units} & \multicolumn{3}{c}{5\% aberrant} & \multicolumn{3}{c}{10\% aberrant} \\
\cmidrule(lr){3-5}\cmidrule(lr){6-8}\cmidrule(lr){9-11}
$(N,T)$ & Estimator & Bias & RMSE & Cover & Bias & RMSE & Cover & Bias & RMSE & Cover \\
\midrule
$(100,50)$ & MG & 0.001 & 0.031 & 94.5 & 0.152 & 0.155 & 36.3 & 0.300 & 0.302 & 0.0 \\
 & Trimmed MG (10\%) & 0.002 & 0.031 & 95.1 & 0.031 & 0.044 & 87.8 & 0.065 & 0.073 & 63.7 \\
 & Median & 0.001 & 0.041 & 95.3 & 0.026 & 0.050 & 91.7 & 0.049 & 0.066 & 82.0 \\
 & MG excl.\ flagged & 0.002 & 0.031 & 94.3 & 0.003 & 0.032 & 95.2 & 0.000 & 0.034 & 93.1 \\
 & \emph{Flagging: TPR / FPR} & \multicolumn{3}{c}{-- / 0.3} & \multicolumn{3}{c}{100.0 / 0.2} & \multicolumn{3}{c}{100.0 / 0.1} \\
\addlinespace
$(200,100)$ & MG & 0.000 & 0.019 & 94.2 & 0.150 & 0.151 & 0.0 & 0.300 & 0.300 & 0.0 \\
 & Trimmed MG (10\%) & 0.000 & 0.021 & 94.3 & 0.023 & 0.031 & 81.6 & 0.054 & 0.058 & 40.4 \\
 & Median & $-$0.001 & 0.029 & 94.8 & 0.020 & 0.035 & 89.9 & 0.044 & 0.054 & 70.4 \\
 & MG excl.\ flagged & 0.000 & 0.019 & 94.2 & 0.000 & 0.019 & 94.0 & 0.000 & 0.020 & 94.4 \\
 & \emph{Flagging: TPR / FPR} & \multicolumn{3}{c}{-- / 0.0} & \multicolumn{3}{c}{100.0 / 0.0} & \multicolumn{3}{c}{100.0 / 0.0} \\
\bottomrule
\end{tabular}
}
\caption{Heterogeneous slopes with aberrant units (slopes shifted by $+3$):
bias, RMSE, and coverage (\%) of nominal 95\% intervals for the centre of
the regular units, $\beta_1=1$, and true / false positive rates (\%) of the
outlier flags. 1{,}000 replications.}
\label{tab:mc3}
\end{table}

Without aberrant units all four estimators are centred on the truth and
have coverage between 94.2\% and 95.3\%; the median pays an efficiency price
(RMSE 0.041 against 0.031 for MG when $N=100$), and the flags mark at most
0.3\% of the regular units. With 5\% aberrant units the MG estimator is
biased by 0.15, and its coverage falls to 36\% for $N=100$ and to zero for
$N=200$: a larger sample only makes the wrong answer more precise. The
trimmed mean and the median remove most of the bias but not all of it,
because the symmetric trimming also cuts regular units from the lower tail,
so their coverage remains below nominal (82\%--92\%) and deteriorates when
10\% of the units are aberrant. Excluding the flagged units restores both
the centre and the coverage (93\%--95\%) in all designs.

With a shift of $+3$ the aberrant units are far outside the distribution of
the regular slopes, which makes detection easy: every aberrant unit is
flagged in every design. To see how the tools behave when detection is hard, we repeated the
$(N,T)=(200,100)$ design with a shift of $+1$, only 2.5 times the
half-width of the distribution of regular slopes (1{,}000 replications;
\code{monte\_carlo.py}, experiment \code{3b}). The flags then identify
55.5\% of the aberrant units when 5\% are aberrant and 39.6\% when 10\% are,
with no false positives. Excluding the flagged units still halves the bias
of the MG estimator with 5\% aberrant units (0.020 against 0.050) but no
longer restores coverage (84.6\%), and with 10\% aberrant units it is no
better than the trimmed mean (bias 0.056 and 0.053). The median has the
smallest bias in both cases (0.022 and 0.044). The diagnostics are
therefore most useful when aberrant units are clearly separated, which is
the case they are designed to reveal; when contamination blends into the
regular distribution, robust summaries and exclusion give similar,
partially corrected answers, and the gap between them and the MG estimate
is itself the informative signal.


\subsection{Unbalanced panels}

Table~\ref{tab:mc4} considers $N=T=100$ with two kinds of missing data.
With \emph{staggered entry}, 30\% of the units enter at a random date
during the first half of the sample and, consistently with
Section~\ref{sec:unbal}, do not interact with other units before they
enter. The experimental unbalanced estimator is compared with the two
balanced alternatives a user of existing software would have: dropping the
late units (and re-standardizing the weights matrix of the remaining units)
or keeping all units but only the periods in which all of them are
observed. With \emph{sporadic} gaps, 2\% of the unit-period cells are
missing at random; the only balanced alternative is to keep the units
without gaps.

\begin{table}[t!]
\centering\scriptsize
{\setlength{\tabcolsep}{4pt}\begin{tabular}{llrrrrrrrrrr}
\toprule
 & & & \multicolumn{3}{c}{$\psi_0$} & \multicolumn{3}{c}{$\rho_1$} & \multicolumn{3}{c}{$\beta_1$} \\
\cmidrule(lr){4-6}\cmidrule(lr){7-9}\cmidrule(lr){10-12}
Design & Estimator & $NT$ & Bias & RMSE & Size & Bias & RMSE & Size & Bias & RMSE & Size \\
\midrule
Staggered & Unbalanced, all units & 9149 & 0.001 & 0.015 & 6.6 & 0.001 & 0.009 & 8.1 & 0.003 & 0.013 & 8.6 \\
 & Drop late units & 6930 & $-$0.070 & 0.073 & 97.6 & 0.003 & 0.012 & 10.0 & 0.004 & 0.016 & 9.5 \\
 & Common period & 5112 & 0.002 & 0.023 & 7.9 & $-$0.002 & 0.013 & 8.9 & 0.002 & 0.018 & 8.2 \\
\addlinespace
Sporadic & Unbalanced, all units & 9507 & $-$0.003 & 0.016 & 7.4 & 0.000 & 0.009 & 9.3 & 0.001 & 0.012 & 7.5 \\
 & Complete units only & 1307 & $-$0.222 & 0.231 & 97.1 & 0.009 & 0.037 & 25.2 & 0.019 & 0.055 & 26.8 \\
\bottomrule
\end{tabular}
}
\caption{Unbalanced panels, $N=T=100$: average number of observations used,
bias, RMSE, and size (\%) of 5\% $t$ tests. 1{,}000 replications.}
\label{tab:mc4}
\end{table}

Under staggered entry the unbalanced estimator is centred on the true
values, and its sizes (6.6\%--8.6\%) are close to those of the 2SIV
estimator in balanced panels of the same dimensions in
Table~\ref{tab:mc1}. Dropping the late units is not a harmless loss of
information: it removes nodes from the network and biases the spatial
coefficient by $-0.07$, so that $t$ tests on $\psi_0$ reject in 98\% of the
replications. Restricting the sample to the common period is valid but
discards almost half of the observations, which raises the RMSE of $\psi_0$
from 0.015 to 0.023. With 2\% sporadic gaps the unbalanced estimator is again
nearly unbiased (bias of $\psi_0$ equal to $-0.003$), whereas keeping only
the units without gaps leaves about 14\% of the observations and severely
biases the spatial coefficient. These results support the usefulness of the
option in designs like these. They do not replace a formal theory, and the
sporadic case in particular could behave worse when missingness is more
frequent or related to the outcome; the option therefore remains labelled
experimental.

